% Bilaga A: metodkapitlet, skrivet för studenterna.  Det står utan
% \chapterprecis: den hör till kapitel som redovisar resultat av en egen
% sökning.  Bilaga B är ett sådant kapitel.  Sökningen gjordes för
% föreläsningen Upprepningar och kapitlet flyttades hit 2026-09-22,
% eftersom den här övningen är det första materialet i kursens läsordning
% som gör påståendet (modules/CLAUDE.md, "Reading order").  Övriga källor
% är belagda i kursens föreläsningar, och de påståenden som saknar belägg
% står i Fortsatt arbete.

\chapter{Att kontrollera ett påstående}
\label{app:verifiera}

Den här övningen gör, som allt undervisningsmaterial, en rad påståenden.
De flesta går att pröva på plats: skriv programmet, kör det och se efter.
Men några gör det inte.  Att det lönar sig att försöka själv innan man
läser ett lösningsförslag, att en funktion med ett enda ansvar är lättare
att ändra, att namn byggda av riktiga ord är lättare att läsa --- det går
inte att se i en terminal.  Sådana påståenden måste kontrolleras mot
källor, och den här bilagan beskriver hur.

Samma fråga möter dig varje gång du läser en lärobok, en webbsida eller
ett svar från en språkmodell, och varje gång du själv ska skriva en
rapport eller granska en kurskamrats kod.  Det är därför arbetsgången
står här och inte bara i lärarens anteckningar.

Ett av övningens påståenden --- att det lönar sig att försöka själv
innan man läser svaret --- prövas här, i en egen sökning:
\cref{app:produktivt} är arbetsgången genomförd, och går att läsa
fristående från resten.  Övriga påståenden är antingen en primärkälla som
gått att slå upp, eller redan kontrollerade i en av kursens föreläsningar
--- och då citeras samma källa, med samma anteckning om hur den hittades,
i stället för att sökas igen.  \Cref{tab:paastaenden} redovisar var vart
och ett av övningens påståenden är kontrollerat, och
\cref{sec:tut-var-fortsatt} de som inte är det.

\begin{table}[!htb]
  \caption{Övningens påståenden som inte går att pröva vid tangentbordet,
    och var vart och ett är kontrollerat.  Kolumnen \enquote{Var det
    står} pekar på avsnittet i övningen som gör påståendet; kolumnen
    \enquote{Kontrollerat} anger vilken bilaga som redovisar sökningen
    --- här eller i en föreläsning ---, eller \enquote{primärkälla} när
    påståendet gäller vad ett dokument säger och dokumentet därför avgör
    saken.  De påståenden som saknar belägg står inte här utan i
    \cref{sec:tut-var-fortsatt}.}
  \label{tab:paastaenden}
  \centering\footnotesize
  \begin{tabular}{@{}p{0.33\linewidth}p{0.18\linewidth}p{0.39\linewidth}@{}}
    \toprule
    Påstående & Var det står & Kontrollerat \\
    \midrule
    Att försöka lösa ett problem innan man får svaret ger bättre
      förståelse, men bara om man sedan jämför sitt försök med svaret
      & Inledningen
      & \Cref{app:produktivt} i den här övningen: måttlig effekt, och
        jämförelsen är ett villkor för den \\
    Varje kunskap i ett program ska ha ett enda ställe (DRY)
      & \Cref{sec:dansa,sec:kakel}
      & Föreläsningen \emph{Algoritmiskt tänkande}, bilaga E: både
        principens upphov och dess förbehåll \\
    En funktion ska aldrig ha mer än ett skäl att ändras (SRP)
      & \Cref{sec:granska}
      & Föreläsningen \emph{Funktioner}, bilaga B: etablerad, men vad
        \emph{ett} ansvar är förblir en bedömningsfråga \\
    Enkelhet är ett designvärde (KISS), men slagordet kommer inte
      bevisligen från Kelly Johnson
      & \Cref{sec:granska}
      & Föreläsningen \emph{Funktioner}, bilaga C: etableringen håller,
        upphovsuppgiften inte \\
    Namn byggda av riktiga ord är lättare att läsa
      & \Cref{sec:granska}
      & Föreläsningen \emph{Variabler och utskrifter}, bilaga B: stöd
        för enskilda regler, inte för stilguiden som helhet \\
    PEP~8 föreskriver versaler för konstanter
      & \Cref{sec:kakel,sec:bmi}
      & Primärkälla, hämtad och läst \\
    PEP~257 är konventionen för docstrings
      & \Cref{sec:granska}
      & Primärkälla, hämtad och läst.  Den flerradiga formen som
        \cref{sec:kakel} använder är däremot inte citerad i
        anteckningen, och står därför i texten som en form koden följer
        --- inte som ett påstående om vad PEP~257 föreskriver \\
    \bottomrule
  \end{tabular}
\end{table}

\section{Gör påståendet till en fråga}
\label{sec:tut-var-fraga}

Ett påstående går att kontrollera först när det är tydligt vad som skulle
kunna visa att det är \emph{fel}.  Formulera det därför som en fråga med
ett svar.  \enquote{Föreskriver PEP~8 versaler för konstanter?} kan
besvaras med ja eller nej; \enquote{docstrings är viktiga} kan inte det,
eftersom \enquote{viktig} är ett omdöme och ingen tänkbar källa skulle
visa att det var fel.  Den som säger så menar oftast något annat --- till
exempel \enquote{slipper den som ska använda en funktion läsa dess kropp
om den har en docstring?} --- och det är en fråga som kan få svaret nej,
och därför kan prövas.

Ofta döljer sig flera påståenden i ett.  Att namnge en princip är två
påståenden: vem den kommer från, och att den är etablerad.  De kräver
olika slags källor, och de kan få olika svar --- KISS i
\cref{sec:granska} är kursens tydligaste exempel på att det ena håller
och det andra inte.

\section{Välj rätt sorts källa}
\label{sec:tut-var-kalla}

Olika frågor besvaras av olika källor.
\begin{description}
  \item[Primärkällan] för vad ett dokument \emph{säger}.  Vill man veta
    vad PEP~8 föreskriver läser man PEP~8, inte en blogg som återger den.
    En text på webben kan ändras, så hämtningsdatum är en del av
    hänvisningen.
  \item[Enskilda studier] för ett specifikt resultat --- att en viss
    missuppfattning har observerats, hos vilka och hur många.  En enskild
    studie visar att något observerats en gång.
  \item[Översikter] (\foreignlanguage{english}{reviews}) och
    metaanalyser för om något är \emph{utbrett} eller vad forskningen
    \emph{sammantaget} säger.  En metaanalys väger samman de publicerade
    försöken och väger därför tyngre än ett av dem.
  \item[Uppslagsverk] för vad ett ord \emph{betyder}.
\end{description}
Var man hittar dem: universitetsbibliotekets söktjänst, och databaser som
Scopus, Web of Science, IEEE Xplore, ACM Digital Library, DBLP (datalogi)
och OpenAlex (öppen; \enquote{öppen} står här för öppna data, inte för
\foreignlanguage{english}{open access}).  Google Scholar täcker brett men
rankar ogenomskinligt.

\section{Sök så att du kan ha fel}
\label{sec:tut-var-sok}

Den som söker på namn hon redan känner till hittar det hon väntar sig.
Sök därför \emph{ämnesdrivet} --- på begrepp, inte författare --- och ställ
samma frågor i flera databaser.  Sök också efter \emph{invändningar}: lägg
till ord som \enquote{critique}, \enquote{limitations} eller
\enquote{no effect}.  Ett påstående som överlevt en motsökning är värt
mer än ett som aldrig prövats, och en invändning som hittas blir ett
\emph{förbehåll} i svaret, inte ett nederlag.  Förbehållet om SRP i
\cref{tab:paastaenden} kom just av en motsökning.  Ett förbehåll kan
också dyka upp utan att någon letat efter det: att vinsten av att
försöka först kräver att man jämför efteråt visade sig i en studie som
stödsökningen hittade (\cref{sec:pf-resultat}), och den räknas inte
mindre för det.

Somliga källor går bara att nå som \emph{kända dokument}: man vet att de
finns och slår upp dem direkt.  Så är det med PEP~8 och PEP~257, som inte
är forskningsartiklar utan gemenskapens egna dokument, och de hämtas
därför från Pythons egen PEP-webbplats.  Det är i sin ordning, bara det
sägs.

\section{Läs i fulltext, och skriv ner hur}
\label{sec:tut-var-fulltext}

En träff i en databas är inte ett belägg.  Sammanfattningen
(\foreignlanguage{english}{abstract}) räcker inte heller: den säger vad
författarna vill ha sagt, inte vad de visat.  Läs hela texten, leta upp
\emph{det ställe} som styrker påståendet och skriv av det ordagrant, med
sidnummer.  Hittar du inget sådant ställe stöder källan inte påståendet,
hur passande titeln än är.  När fulltexten inte går att nå ska det stå i
anteckningen, och påståendet inte sägas starkare än sammanfattningen
tillåter.  Så är det för flera av den här övningens källor --- bland dem
de som står på motsidan i \cref{app:produktivt} --- och för en annan
gick fulltexten att läsa en gång men inte att nå en andra gång; båda
sakerna står i anteckningen, ordagrant.

Kravet gäller också --- kanske särskilt --- källor du har ärvt.  En
hänvisning som redan står i ett dokument är inte kontrollerad förrän
\emph{du} har läst stället den pekar på.  Den här övningen ärver de
flesta av sina källor från kursens föreläsningar, och det som gör det
försvarbart är att anteckningen följer med: i litteraturlistans källfil
står, ovanför varje post, dels den ursprungliga anteckningen från den
föreläsning som läste källan, dels en anteckning om att den är kopierad
hit och inte läst om.  Källorna i \cref{app:produktivt} har i stället
sin egen anteckning.  Fälten är fasta:
\begin{description}
  \item[\texttt{CLAIM}] vilket påstående källan ska styrka;
  \item[\texttt{FOUND-VIA}] hur den hittades --- sökfråga och databas,
    eller \enquote{känt dokument};
  \item[\texttt{PICKED}] varför just den, framför andra träffar;
  \item[\texttt{QUOTE}] det ordagranna citatet, med sida;
  \item[\texttt{VERIFIED}] att fulltexten lästs, och var;
  \item[\texttt{COUNTER}] vad motsökningen gav;
  \item[\texttt{DATE}] när.
\end{description}
Anteckningen ser byråkratisk ut och tar ett par minuter per källa.  Den
är också det enda som gör det möjligt att om ett år svara på frågan
\enquote{hur vet du det?}

\section{Dra slutsatsen --- med förbehåll}
\label{sec:tut-var-slutsats}

Svaret skrivs ut tillsammans med det som talar emot.  Inget av
påståendena i \cref{tab:paastaenden} är obetingat, och texten i övningen
säger dem därför inte starkare än bilagorna tillåter: DRY och SRP är
tumregler med kända förbehåll, PEP~8 stöds för enskilda regler och inte
som helhet, och vinsten av att försöka först försvinner om man hoppar
över jämförelsen efteråt.  När du själv motiverar ett val i din kod ---
till exempel i granskningen av en kurskamrats lösning --- duger samma
måttstock: säg vad du vet, hur du vet det, och vad du inte vet.

\section{Fortsatt arbete}
\label{sec:tut-var-fortsatt}

Några påståenden i övningen är \emph{inte} kontrollerade, och de står
här i stället för att döljas.  Alla är påståenden om \emph{nytta}, och
sådana är lätta att skriva utan att tänka på att de går att pröva.
\begin{description}
  \item[Återanvändbarhet] Att en funktion som lämnar tillbaka sitt
    resultat går att använda till mer än en som skriver ut det
    (\cref{sec:diska}).  Att den \emph{kan} användas på fler sätt följer
    av vad ett returvärde är; att det gör någon skillnad för den som
    skriver eller underhåller programmet är en empirisk fråga som ingen
    av kursens bilagor har prövat.
  \item[Docstringens nytta] Att en docstring gör att den som ska använda
    funktionen slipper läsa dess kropp (\cref{sec:granska}
    och lärandemål~4).
    Föreläsningen \emph{Funktioner} håller sig medvetet till det som går
    att visa --- att \mintinline{python}{help} och \texttt{pydoc} läser
    docstringen --- och påstår ingenting om läsaren.  Den här övningen
    går ett steg längre, och steget är inte belagt.
  \item[Verb för funktioner] Att en funktion bör heta något som säger
    vad den gör och en variabel något som säger vad den innehåller
    (\cref{sec:granska}).  PEP~8 föreskriver inte ordklassen, och
    sökningen bakom bilaga B i föreläsningen \emph{Variabler och
    utskrifter} prövade om namn byggda av riktiga ord är lättare att
    läsa --- inte om verb är bättre än substantiv.
  \item[Att granska andras kod] Att läsa och granska någon annans
    program gör den egna koden tydligare (\cref{sec:granska}).  Det är
    kursens eget motiv för momentet och ett av dess lärandemål, men
    ingen av kursens bilagor har sökt efter underlag för det, och
    kamratgranskning är ett välstuderat område där ett svar borde gå att
    hitta.
\end{description}
Var och en av dem är en fråga att ställa i en egen sökning, tvåsidigt,
enligt arbetsgången i \cref{sec:tut-var-fraga,sec:tut-var-kalla,%
sec:tut-var-sok,sec:tut-var-fulltext,sec:tut-var-slutsats}.  Tills det är gjort står de som de står här:
som skäl författaren tror på, inte som belägg.

\section{En anmärkning om språkmodeller}
\label{sec:tut-var-sprakmodeller}

Sökningen i \cref{app:produktivt} gav hundratals träffar.  En
språkmodell användes för att \emph{sortera} dem --- hör träffen till
ämnet över huvud taget, och åt vilket håll pekar den? --- så att
gallringen kunde granskas.  Så har det gått till också i de sökningar
som föreläsningarna \emph{Algoritmiskt tänkande}, \emph{Funktioner} och
\emph{Variabler och utskrifter} redovisar.  Lägg märke till
arbetsfördelningen: modellen sorterade, men varje källa som citeras ska
läsas och väljas av en människa, och även sorteringen ska bekräftas av en
människa.  Modellens klassning står därför med sin konfidens i
\cref{tab:sok-pf}, så att den kan ifrågasättas.  Använd gärna
språkmodeller för att hitta och sortera; låt dem aldrig vara det sista
ledet.  De kan göra fel, men det är du själv som måste stå till svars
för innehållet.

\chapter{Lönar det sig att låta studenten försöka först?}
\label{app:produktivt}
\chapterprecis{Författaren har ännu inte granskat resultaten i den här
  bilagan i sin helhet.}

Den här övningen är byggd så att du ska försöka först.  Varje uppgift
står före sitt lösningsförslag --- vad skriver programmet ut, hur skulle du
dela upp det, vad skulle du göra annorlunda? --- och kursens föreläsningar
är byggda på samma sätt: före varje förklaring står en fråga, och svaret
kommer efteråt.  Det är ett medvetet val, och ett som går att
ifrågasätta: det vore lika rimligt att först få veta hur något fungerar och sedan öva.  Frågan är alltså:
\emph{leder det till bättre lärande att låta studenten försöka lösa ett
problem innan begreppet lärs ut, än att lära ut begreppet först?}

Frågan har egna namn i forskningen ---
\emph{\foreignlanguage{english}{productive failure}} och
\emph{\foreignlanguage{english}{problem solving prior to instruction}},
förkortat PS-I --- och den är omdiskuterad, vilket gör den till ett bra
exempel på hur en motsökning ska gå till.  \Cref{sec:pf-metod} beskriver
sökningen, \cref{sec:pf-resultat} vad den fann, \cref{sec:pf-diskussion}
vad som begränsar svaret och \cref{sec:pf-slutsats} ger det;
\cref{sec:pf-fortsatt} säger vad som återstår att undersöka, och
\cref{tab:sok-pf} sist i kapitlet listar samtliga träffar som rör
påståendet.

\section{Metod}
\label{sec:pf-metod}

Sökningarna gjordes 2026-09-03 och 2026-09-04 med verktyget
\texttt{scholar} mot fyra databaser: OpenAlex (OA; inte \foreignlanguage{english}{open access}), DBLP,
Web of Science (WoS) och Scopus.  Två databaser som annars ingår saknas:
Semantic Scholar avvisade sin nyckel, och IEEE Xplore hade slut på dagens
kvot.  En femte, arXiv, prövades på de två första frågorna och gav 40
träffar var --- samtliga om maskininlärning, matchade på ordet
\foreignlanguage{english}{learning}.  arXiv indexerar fysik och datalogi,
inte pedagogisk psykologi, och togs därför ur resten av sökningen; det sägs
här därför att en utelämnad databas alltid ska motiveras.

Nio ämnesfrågor ställdes, fyra som söker stöd och fem som söker
invändningar, och varje fråga ställdes till alla fyra databaserna med samma
nyckelord.
Frågorna är ämnesdrivna: ingen av dem nämner en författare, trots att
området domineras av ett fåtal namn.  Motfrågorna använder dels samma
begrepp som stödfrågorna med tillägg som
\foreignlanguage{english}{critique} och
\foreignlanguage{english}{limitations}, dels den motsatta ståndpunktens egna
termer --- \foreignlanguage{english}{minimal guidance},
\foreignlanguage{english}{explicit instruction},
\foreignlanguage{english}{element interactivity} --- eftersom en motsökning
som bara söker på det man tror på inte är någon motsökning.
Till dem kommer en tionde sökning, K1, som inte söker efter ämnet utan
efter ett känt dokument.  Den behövdes: ämnesfrågorna gav bara författarnas
egen preprint av metaanalysen, aldrig den publicerade artikeln, och en
preprint duger inte som belägg (\cref{app:verifiera}).  Artikeln hämtades
därför på sin titel dagen därpå, och sökningen står i tabellen som
varje annan.
\Cref{tab:pf-queries} visar frågorna och antalet träffar per databas.

\begin{table}[htbp]
  \centering
  \caption{Sökfrågor och antal träffar per databas, 2026-09-03
    (S1--M5) och 2026-09-04 (K1), verktyget \texttt{scholar}.
    S = stödfråga, M = motfråga, K = känt dokument.  Databaser:
    OpenAlex (OA; inte \foreignlanguage{english}{open access}), DBLP, Web of
    Science (WoS), Scopus.  Högst 40 träffar hämtades per databas och fråga,
    så 40 betyder \enquote{minst 40}.  Web of Science fick nyckelorden som
    \texttt{TS=(\dots)} och Scopus som \texttt{TITLE-ABS-KEY(\dots)}, medan
    OpenAlex fick boolesk fritext och DBLP en ordlista där varje ord måste
    matcha --- vilket är varför DBLP ger noll på nästan allt: det är en
    datalogidatabas, och frågorna är pedagogiska.}
  \label{tab:pf-queries}
  \footnotesize
  \begin{tabular}{@{}lp{0.50\linewidth}rrrr@{}}
    \toprule
    & Nyckelord & OA & DBLP & WoS & Scopus \\
    \midrule
    S1 & \enquote{productive failure} learning outcomes
       & 40 & 0 & 40 & 40 \\
    S2 & \enquote{problem solving prior to instruction} learning
       & 40 & 1 & 26 & 40 \\
    S3 & \enquote{invention activities} \enquote{preparation for future
         learning}
       & 40 & 0 & 4 & 40 \\
    S4 & exploration before instruction conceptual understanding
       & 39 & 0 & 11 & 40 \\
    \midrule
    M1 & \enquote{productive failure} (critique, limitations eller
         replication)
       & 40 & 0 & 16 & 40 \\
    M2 & \enquote{problem solving prior to instruction}
         (\enquote{no difference}, \enquote{no effect} eller null)
       & 40 & 0 & 1 & 20 \\
    M3 & \enquote{minimal guidance} instruction novices worked examples
       & 40 & 0 & 2 & 40 \\
    M4 & \enquote{explicit instruction} versus \enquote{discovery
         learning} achievement
       & 40 & 0 & 0 & 40 \\
    M5 & \enquote{productive failure} boundary conditions prior knowledge
       & 40 & 0 & 1 & 40 \\
    \midrule
    K1 & \enquote{When Problem Solving Followed by Instruction Works}
         (känt dokument; högst 10 träffar per databas)
       & 10 & 0 & 1 & 10 \\
    \bottomrule
  \end{tabular}
\end{table}

% Author's note (verification and reproduction): see the block above the
% hit-list \input at the end of this chapter.
Sammanlagt gav sökningarna 536 poster.  Fyra av dem är dubbletter --- samma
arbete hämtat ur flera databaser, och i ett fall en preprint tillsammans med
den publicerade artikeln --- så antalet unika arbeten är 532, och det är det
talet \cref{tab:sok-pf} räknar med.  Alla klassades: först med en
språkmodell mot en beskrivning av frågan och fyra kategorier --- stöder
påståendet, kvalificerar eller motsäger det, angränsande delämne, annat ämne
--- och sedan för hand.  Varje post som modellen satt i den andra kategorin
lästes, liksom varje post där modellens konfidens låg under 0,7, och tio
beslut fattades för hand.  Fyra av dem ändrade modellens: tre poster som
modellen behållit visade sig vara systembeskrivningar eller protokoll utan
resultat, och en som den kastat --- en översikt över när lärande utan
föregående undervisning fungerar --- hörde tvärtom till kärnan av frågan.
Tre poster flyttades dessutom mellan kategorier; en studie som modellen läst
som en invändning rapporterar i själva verket stöd, med ett förbehåll om
åldersgrupp.

Vid genomläsningen av den färdiga tabellen rättades fyra saker för hand, och
de står här därför att en rättelse som inte redovisas är en rättelse man
inte kan kontrollera.  Tre rader hade modellen märkt både med
\texttt{qualifies\_claim} och med ett ämnesord som säger att posten inte
gäller frågan (\texttt{adjacent\_subtopic}, \texttt{not\_ps\_i\_specific});
regeln ovan säger att sådana ämnesord ska göra posten angränsande, så de
tre flyttades dit.  Preprinten och artikeln av javelinstudien är ett och
samma arbete och står nu på en rad.  Preprinten av metaanalysen fick en egen
etikett, så att den inte läses som den citerade källan.  Och
leverantörskolumnen redovisar ämnesfrågorna: de tre rader som K1 råkade
skriva om ställdes tillbaka.  Kontrollen av de två raderna
\enquote{Motivation and transfer} gav däremot ingen rättelse --- det är två
skilda publikationer, ett konferensbidrag från 2009 utan \textsc{doi} och en
tidskriftsartikel från 2012, med olika sammanfattningar --- så de står kvar
som två rader.

En sak i \cref{tab:sok-pf} ska läsas med försiktighet.  Modellen fick
beskriva varje träff med egna ämnesord, och 81 poster fick ämnesord men
ingen riktning --- det är oftast studier av produktivt misslyckande där
modellen inte skrev ut åt vilket håll resultatet pekade.  De räknas här som
angränsande.  Det underskattar stödsidan snarare än överskattar den, vilket
är åt rätt håll när man prövar sitt eget påstående.

\section{Resultat}
\label{sec:pf-resultat}

\paragraph{Effekten finns, och den är måttlig.}  Den enskilt tydligaste
studien kommer från \textcite[1008]{Kapur2014}: två randomiserade försök
där eleverna antingen fick lösa ett problem först och undervisas sedan,
eller tvärtom.
\enquote{Two randomized-controlled studies found that
both methods lead to high levels of procedural knowledge. However, students
who engaged in problem solving before being taught demonstrated
significantly greater conceptual understanding and ability to transfer to
novel problems than those who were taught first}.  Lägg märke till att
skillnaden inte ligger i \emph{procedurkunskapen} --- att kunna utföra
stegen --- utan i förståelsen och i förmågan att använda den på nya problem.

Att det inte bara är en studie visar metaanalysen av
\textcite{SinhaKapur2021}, som väger samman de publicerade försöken:
\textcquote[761]{SinhaKapur2021}{a significant, moderate effect in favor of
PS-I (Hedge's g 0.36 [95\% confidence interval 0.20; 0.51])}, och starkare
när upplägget följer produktivt misslyckande troget.  En metaanalys väger
tyngre än en enskild studie (\cref{app:verifiera}), och det är därför den
här bilagan lutar sig mot den.

\paragraph{Men bara om efterarbetet görs rätt.}  Det viktigaste förbehållet
kommer inte från kritikerna utan från fältet självt.
\textcite{LoiblRummel2013} varierade inte \emph{om} undervisningen sköts
upp, utan \emph{vad} som hände efteråt, och fann att
\textcquote[296]{LoiblRummel2013}{comparing and contrasting typical student
solutions is a prerequisite for the effectiveness of problem-solving prior
to instruction}.  Utan den jämförelsen försvann hela vinsten:
\textcquote[296]{LoiblRummel2013}{Problem-solving prior to instruction was
no more effective than standard instruction prior to problem-solving}.  Att
skjuta upp svaret räcker alltså inte; det är kontrasten mellan det
studenten själv försökte och det riktiga svaret som gör jobbet.

\paragraph{Och inte överallt.}  Motfrågornas träffar sätter gränser.
\textcite{Fyfe2014} fann det omvända --- undervisning först var bättre ---
hos andra- och tredjeklassare med låga förkunskaper i matematik:
\textcquote[2]{Fyfe2014}{Providing conceptual instruction prior to problem
solving was the more effective sequencing of activities than the reverse}.
Samma gräns finns i metaanalysen av PS-I, som noterar att
\textcquote[780]{SinhaKapur2021}{Contrasting trends were, however, observed
for younger age learners (second to fifth graders) and for the learning of
domain-general skills, for which effect sizes favored I-PS}
--- där I-PS är den omvända ordningen, undervisning före problemlösning.
\textcite{Ashman2020} pekar ut en annan gräns, i termer av hur många delar
som måste hållas i huvudet samtidigt:
\textcquote{Ashman2020}{for learning where element interactivity is high,
explicit instruction should precede problem-solving}.  Och
\textcite{Nachtigall2020} försökte upprepa effekten utanför
STEM-ämnena (\foreignlanguage{english}{science, technology, engineering and
mathematics}) och misslyckades:
\textcquote{Nachtigall2020}{the findings did not replicate the PF effect on
learning in a non-STEM domain} --- PF är förkortningen för
\foreignlanguage{english}{productive failure}.  Till det kommer ett
oavgjort resultat:
\textcite{ChaseKlahr2017} fann att
\textcquote{ChaseKlahr2017}{students in both conditions made equally large
gains}, och drar slutsatsen att breda påståenden om vilken ordning som är
bäst \textcquote{ChaseKlahr2017}{should be conditionalized by particular
instructional contexts}.

\paragraph{Den principiella invändningen.}  Den starkaste finns i
\textcite{KirschnerSwellerClark2006}, som sammanfattar ett halvsekels
forskning så här:
\textcquote[75]{KirschnerSwellerClark2006}{minimally guided instruction is
less effective and less efficient than instructional approaches that place a
strong emphasis on guidance of the student learning process. The advantage
of guidance begins to recede only when learners have sufficiently high prior
knowledge to provide \mkbibquote{internal} guidance}.  Deras måltavla är
undervisning som lämnar studenten att upptäcka själv utan hjälp.  Det är
inte vad kursens material gör --- frågorna och uppgifterna är korta,
konkreta och följs av svaret --- men invändningen sätter fingret på var
gränsen går.

\paragraph{Prövat i just det här ämnet.}  En enda studie bland
träffarna gör försöket i programmeringsundervisning.
\textcite{Suriyaarachchi2025} lät nybörjare i Python antingen brottas med
ett problem först eller få genomgången först, och fann att
\textcquote{Suriyaarachchi2025}{although there was no difference in initial
learning outcomes between the groups, students who followed the PF approach
showed better knowledge retention and performance on delayed but similar
tasks}.  Ingen omedelbar vinst alltså, men en vinst två veckor senare ---
och studien är liten: tjugo deltagare totalt.

\needspace{6\baselineskip}
\section{Diskussion och begränsningar}
\label{sec:pf-diskussion}

Flera saker begränsar svaret.

Semantic Scholar och IEEE Xplore saknas, och arXiv togs bort med angivet
skäl, så täckningen vilar på fyra databaser.  Träffantalen i
\cref{tab:pf-queries} är tak, inte totaler: 40 betyder \enquote{minst 40},
och för de flesta frågorna var det verkliga antalet större.

Tre av de citerade källorna --- \textcite{Ashman2020},
\textcite{Nachtigall2020} och \textcite{ChaseKlahr2017} --- är bara lästa i
sammanfattning, eftersom förlagen inte ger fulltexten och ingen öppen kopia
hittades.  De citeras därför bara för det deras egna sammanfattningar säger,
och inget mer.  Det gäller alla tre på \emph{motsidan}, vilket är värt att
notera: den sida som är svårast att läsa i fulltext är den som är lättast
att avfärda.

En fjärde källa gick inte att läsa alls.  Motsökningen fann
\enquote{Reverse the routine: Problem solving before instruction improves
conceptual knowledge in undergraduate physics} (2018), som efter sin titel
är den närmaste replikationen på universitetsnivå.  Förlaget ger varken
fulltext eller en öppen kopia, och den citeras därför inte --- den står i
\cref{tab:sok-pf} som noterad men inte citerad.  Reglerna gäller åt båda
hållen: en källa man inte har läst räknas inte, hur väl den än passar.

Slutligen är den enda studien i programmeringsämnet liten, och dess vinst
syns bara i eftertestet.  Att generalisera från tjugo deltagare vore
oförsiktigt; den väger här som en indikation, inte som ett bevis.

\section{Slutsats}
\label{sec:pf-slutsats}

Ja, med två förbehåll som båda styr hur kursens material är skrivet.
Effekten finns och är måttlig i genomsnitt över de publicerade försöken, och
den ligger i förståelse och överföring snarare än i att kunna utföra stegen.
Men den uppstår inte av att svaret skjuts upp.  Den uppstår av att
studentens eget försök sedan ställs mot det riktiga svaret --- utan den
jämförelsen är ordningen likgiltig.  Därför följs varje uppgift i den här
övningen av ett lösningsförslag, och därför ber övningen dig jämföra ditt
försök med förslaget i stället för att bara läsa det.  I kursens
föreläsningar följs varje fråga på samma sätt av svaret \emph{och} av en
uttalad kontrast.

Det andra förbehållet är var effekten gäller.  Den vänder för barn i
lågstadieåldern och för allmänna färdigheter, och den uteblir när materialet
kräver att många delar hålls i huvudet samtidigt.  Kursens studenter är
vuxna, ämnet är konkret, och programmen är små: föreläsningarnas exempel är
ofta bara några rader långa, och övningens uppgifter delar sina program i
funktioner på några rader var --- vi ligger alltså, så långt det går att
bedöma, inom det område där stödet finns, men det är värt att veta att
gränsen finns och var den går.  Att invändningarna är
starkast just där de är svårast att läsa i fulltext är den sista lärdomen:
ett påstående är inte kontrollerat förrän motsidan har fått samma chans.

\section{Fortsatt arbete}
\label{sec:pf-fortsatt}

\Cref{sec:pf-diskussion} räknar upp vad som begränsar svaret; det här
avsnittet säger vad som ska göras åt det.  Frågorna ska ställas en gång
till på Semantic Scholar och IEEE Xplore när de svarar, och träfftaket
höjas för de frågor där fyrtio betydde \enquote{minst fyrtio}, så att
antalen blir totaler och inte tak.  De tre källor som bara är lästa i
sammanfattning ska läsas i
fulltext\autocite{Ashman2020,Nachtigall2020,ChaseKlahr2017}; att alla
tre står på motsidan gör dem viktigare, inte mindre viktiga.  Och Weaver
med fleras artikel från 2018 i \emph{Contemporary Educational
Psychology}, som efter sin titel är den närmaste replikationen på
universitetsnivå, ska skaffas genom ett bibliotek: den skulle pröva om
effekten håller i den åldersgrupp kursen faktiskt
undervisar.\footnote{\textsc{doi} 10.1016/j.cedpsych.2017.12.003.}

För ett säkrare svar saknas två saker, och den ena är lätt att namnge:
ämnet.  En enda av de 532 posterna prövar ordningen i
programmeringsundervisning, med tjugo deltagare, och vinsten syns bara i
eftertestet.  Ett större kontrollerat försök i en inledande
programmeringskurs, med ett fördröjt eftertest, är det som skulle göra
svaret till ett svar om just den här sortens undervisning.  Den andra
saken följer av kapitlets egen slutsats: om effekten kommer av
jämförelsen mellan det egna försöket och det riktiga svaret, och inte av
att svaret skjuts upp, så borde jämförelsen prövas som en egen faktor
--- försök först med efterföljande kontrast mot försök först utan.
Ingen studie bland träffarna skiljer de två åt.  Därtill vet ingen var
gränsen för hur mycket som måste hållas i huvudet samtidigt går för så
små program som kursens.  Föll ett sådant försök ut
till ordningens fördel skulle kursens uppläggning vila på ett belägg
ur sitt eget ämne i stället för på ett genomsnitt över andra ämnen; föll
det åt andra hållet vore frågorna fortfarande värda att ställa, men
förklaringen skulle behöva komma först.

\section{Träffar som rör påståendet}
\label{sec:pf-traffar}

\Cref{tab:sok-pf} listar de poster som rör påståendet --- citerade,
stödjande och kvalificerande --- med skälet till varje inkludering; de
angränsande och de felträffar som uteslöts anges som antal i tabelltexten,
och antalet träffar per fråga står i \cref{tab:pf-queries}.
% Author's note (verification and reproduction):
%   moved           : 2026-09-22 from modules/iterations/slides
%                     (Upprepningar, bilaga C there), with its exports and
%                     bib entries, because this tutorial is the first deck
%                     in the reading order that makes the claim.  Session
%                     name and export paths are unchanged.  Wording that
%                     named the lecture now names this tutorial; the
%                     search, counts and conclusions are unchanged.
%   scholar session : iterations-productive-failure
%   exports         : litteratursokning/produktivt-misslyckande.csv and
%                     litteratursokning/produktivt-misslyckande-full.tex
%   searched        : 2026-09-03 (S1-S4, M1-M5); 2026-09-04 (K1);
%                     providers openalex dblp wos scopus
%                     (s2: key rejected HTTP 403; ieee: daily quota
%                     exhausted -- both absent.  arxiv: tried on S1 and S2,
%                     40 hits each, all machine-learning false hits on the
%                     word "learning"; dropped from the remaining queries
%                     with that reason stated in the method section)
%                     limit 40 per provider and query
%   known item (K1) : scholar search
%                     '"When Problem Solving Followed by Instruction Works"'
%                     -p openalex -p dblp -p wos -p scopus -l 10
%                     -n iterations-productive-failure
%                     (openalex 10, dblp 0, wos 1, scopus 10; the Review of
%                     Educational Research article 10.3102/00346543211019105
%                     is the first hit from openalex, wos and scopus).  Run
%                     because S1-S4/M1-M5 returned only the authors' OSF
%                     preprint of the same work (10.35542/osf.io/83p7e).
%                     Four other hits were new to the session and were
%                     decided by hand without being read:
%                     scholar sessions decide iterations-productive-failure
%                       --keep --doi 10.3102/00346543211019105
%                       -t supports-claim -t supports_claim -t meta_analysis
%                       -t preparatory_approaches -t pf_meta
%                     scholar sessions decide iterations-productive-failure
%                       --discard --doi 10.17605/osf.io/mknsd
%                       -t off-topic-false-hit -t other_topic
%                     scholar sessions decide iterations-productive-failure
%                       --discard --doi 10.1038/s41539-023-00197-4
%                       --doi 10.1007/s11423-022-10104-0 --doi 10.1002/sce.70060
%                       -t adjacent-subtopic -t not_ps_i_specific
%                     Adding K1's results merged two previously separate
%                     records of "Delaying Instruction Alone Doesn't Work"
%                     into one, so the session went 532 -> 536 records.
%   classification  : scholar llm context + scholar llm classify
%                     --no-examples --no-enrich, model
%                     github_copilot/gpt-5.4; 10 human decisions via
%                     scholar sessions decide (read: every qualifies-claim
%                     row and every row with confidence < 0.7), of which
%                     4 changed the model's decision and 3 changed its
%                     category
%   tag vocabulary  : the model was allowed free-form tags; they were
%                     mapped onto the audit table's four categories by a
%                     stated rule (supports_claim -> supports-claim,
%                     qualifies_claim -> qualifies-claim, adjacent* and
%                     not_ps_i* -> adjacent-subtopic, other_topic and the
%                     machine-learning tags -> off-topic-false-hit, and
%                     81 rows the model gave topic words but no direction
%                     -> adjacent-subtopic, which understates the
%                     supporting side; said so in the method section)
%   hit-list table  : exported, then hand-corrected (see below).
%                     scholar sessions export iterations-productive-failure
%                     -f table --bearing-only --lang sv --label sok-pf
%                     --track "produktivt misslyckande"
%                     --theme pf_core="kärnexperimentet"
%                     --theme pf_meta="metaanalys av PS-I"
%                     --theme pf_cs1="prövat i programmeringsundervisning"
%                     --theme pf_consolidation="efterarbetet är förutsättningen"
%                     --theme pf_consolidation_dup="efterarbetet är förutsättningen"
%                     --theme pf_guidance="kritiken mot svag vägledning"
%                     --theme pf_guidance_dup="kritiken mot svag vägledning"
%                     --theme pf_reverse="undervisning först var bättre"
%                     --theme pf_element="elementinteraktivitet som gräns"
%                     --theme pf_domain="replikerades inte utanför STEM"
%                     --theme pf_tie="ingen skillnad"
%                     -o litteratursokning/produktivt-misslyckande-full
%   hand edits      : the exported table was edited by hand afterwards;
%                     scripted form kept out of tree, every change stated in
%                     the method section.  (1) The three rows the model
%                     tagged both qualifies_claim and adjacent_subtopic /
%                     not_ps_i_specific (otoscopy simulator; Preschool
%                     Children's Science Learning; Revisiting cognitive load
%                     theory) were moved to the adjacent count, which is
%                     what the stated mapping rule says those tags mean.
%                     (2) The javelin preprint (10.35542/osf.io/rd2ya) and
%                     article (10.1016/j.learninstruc.2023.101828) are one
%                     work and are now one row.  (3) The OSF preprint of the
%                     meta-analysis keeps its row but is labelled as the
%                     preprint, not as a cited source.  (4) K1 restamped its
%                     own provider list on three rows it touched; those
%                     cells were restored to the providers of the topic
%                     searches, and the new row carries the providers that
%                     returned it (OA, WoS, Scopus; DBLP gave 0 hits).
%                     (5) Long snake_case tags got \_\allowbreak{} so the
%                     Skäl column can break them, and that column is set in
%                     \footnotesize.  Counts in the caption follow: 9 cited,
%                     22 supporting (one of them the preprint row), 21
%                     qualifying, 364 adjacent, 116 false hits, 532 unique
%                     works out of 536 records.
%   provenance      : ltnotes.bib, keys Kapur2014 SinhaKapur2021
%                     Suriyaarachchi2025 LoiblRummel2013
%                     KirschnerSwellerClark2006 Fyfe2014 Ashman2020
%                     Nachtigall2020 ChaseKlahr2017 (blocks CLAIM /
%                     FOUND-VIA / PICKED / QUOTE / VERIFIED / COUNTER /
%                     DATE)
%   known items     : SinhaKapur2021 -- the search returned the authors'
%                     OSF preprint (10.35542/osf.io/83p7e); the published
%                     Review of Educational Research article
%                     (10.3102/00346543211019105) was retrieved from it as
%                     a known item and is the version read and cited.
%                     The preprint itself was NOT read.
%   abstract only   : Ashman2020 (ERIC EJ1245274), Nachtigall2020 (ERIC
%                     EJ1278375), ChaseKlahr2017 (publisher page, checked
%                     against ERIC EJ1157912) -- flagged in the discussion
%   not read        : Weaver et al. 2018, 10.1016/j.cedpsych.2017.12.003
%                     (closed access, no repository copy; listed in the
%                     hit table as recorded-not-cited)
% --- scholar LaTeX fragment --------------------------------------------
% Generated by: scholar sessions export -f table
% Session: iterations-productive-failure
% \input this file from the body of a document whose
% preamble loads:
%   \usepackage{longtable}
%   \usepackage{booktabs}
%   \usepackage{array}
% ----------------------------------------------------------------------

\begingroup\footnotesize
\begin{longtable}{@{}%
  >{\raggedright\arraybackslash}p{0.42\textwidth}c%
  >{\raggedright\arraybackslash}p{0.13\textwidth}%
  >{\raggedright\arraybackslash\footnotesize}p{0.24\textwidth}@{}}
\caption{Träffar som rör påståendet, produktivt misslyckande (9 citerade, 22 som stöder påståendet --- en av dem är preprintversionen av den citerade metaanalysen --- och 21 som kvalificerar eller motsäger det, av 532 unika arbeten; de 364 angränsande och 116 felträffarna listas inte). Skälet till varje inkludering står i sista kolumnen; för maskinklassade rader anges modellens konfidens. Databaser: OpenAlex (OA; inte open access), DBLP, Web of Science (WoS), Scopus.}\label{tab:sok-pf}\\
\toprule
Titel & År & Leverantör & Skäl \\
\midrule
\endfirsthead
\multicolumn{4}{@{}l}{\emph{\tablename~\thetable{} (forts.)}}\\
\toprule
Titel & År & Leverantör & Skäl \\
\midrule
\endhead
\midrule \multicolumn{4}{r@{}}{\emph{forts.\ på nästa sida}}\\
\endfoot
\bottomrule
\endlastfoot
An alternative time for telling: When conceptual instruction prior to problem solving improves mathematical knowledge & 2014 & OA & \emph{citerad}: undervisning först var bättre (qualifies\_\allowbreak{}claim; instruction\_\allowbreak{}first\_\allowbreak{}better; math; children) \\
Delaying Instruction Alone Doesn't Work: Comparing and Contrasting Student Solutions is Necessary for Learning from Problem-Solving Prior to Instruction. & 2013 & OA, DBLP & \emph{citerad}: efterarbetet är förutsättningen (qualifies\_\allowbreak{}claim; instruction\_\allowbreak{}phase\_\allowbreak{}matters; productive\_\allowbreak{}failure) \\
From problem-solving to sensemaking: A comparative meta-analysis of preparatory approaches for future learning & 2021 & OA & preprint (10.35542/osf.io/83p7e) av den citerade metaanalysen; inte själv läst eller citerad (meta\_\allowbreak{}analysis; preparatory\_\allowbreak{}approaches; qualifies\_\allowbreak{}claim; supports\_\allowbreak{}claim) \\
Invention Versus Direct Instruction: For Some Content, It's a Tie & 2017 & Scopus & \emph{citerad}: ingen skillnad (qualifies\_\allowbreak{}claim; invention\_\allowbreak{}vs\_\allowbreak{}direct\_\allowbreak{}instruction; tie\_\allowbreak{}results) \\
Investigating the Use of Productive Failure as a Design Paradigm for Learning Introductory Python Programming & 2025 & WoS & \emph{citerad}: prövat i programmeringsundervisning (productive\_\allowbreak{}failure; computer\_\allowbreak{}science; null\_\allowbreak{}or\_\allowbreak{}mixed\_\allowbreak{}results; qualifies\_\allowbreak{}claim) \\
Problem-solving or Explicit Instruction: Which Should Go First When Element Interactivity Is High? & 2020 & Scopus & \emph{citerad}: elementinteraktivitet som gräns (qualifies\_\allowbreak{}claim; cognitive\_\allowbreak{}load\_\allowbreak{}critique; element\_\allowbreak{}interactivity) \\
Productive Failure in Learning Math & 2014 & OA & \emph{citerad}: kärnexperimentet (supports\_\allowbreak{}claim; productive\_\allowbreak{}failure; math; conceptual\_\allowbreak{}understanding; transfer) \\
When failure fails to be productive: probing the effectiveness of productive failure for learning beyond STEM domains & 2020 & OA & \emph{citerad}: replikerades inte utanför STEM (qualifies\_\allowbreak{}claim; productive\_\allowbreak{}failure\_\allowbreak{}fails; non\_\allowbreak{}stem; boundary\_\allowbreak{}conditions) \\
When Problem Solving Followed by Instruction Works: Evidence for Productive Failure & 2021 & OA, WoS, Scopus & \emph{citerad}: metaanalys av PS-I (supports\_\allowbreak{}claim; meta\_\allowbreak{}analysis; preparatory\_\allowbreak{}approaches) \\
Why Minimal Guidance During Instruction Does Not Work: An Analysis of the Failure of Constructivist, Discovery, Problem-Based, Experiential, and Inquiry-Based Teaching & 2006 & OA, WoS & \emph{citerad}: kritiken mot svag vägledning (qualifies\_\allowbreak{}claim; cognitive\_\allowbreak{}load\_\allowbreak{}critique; minimal\_\allowbreak{}guidance) \\
Active before passive tasks improve long-term visual learning in difficult-to-classify skin lesions & 2023 & WoS & stöder påståendet (0.91) (active\_\allowbreak{}before\_\allowbreak{}passive; supports\_\allowbreak{}claim; sequence\_\allowbreak{}comparison; transfer) \\
Benefits of Problem-Solving First for the Success of Cooperative Learning with Videos & 2026 & Scopus & stöder påståendet (0.78) (supports\_\allowbreak{}claim; problem\_\allowbreak{}solving\_\allowbreak{}first; cooperative\_\allowbreak{}learning; videos) \\
Bringing exploratory learning online: problem-solving before instruction improves remote undergraduate physics learning & 2023 & WoS & stöder påståendet (0.98) (exploratory\_\allowbreak{}learning; physics; supports\_\allowbreak{}claim; remote\_\allowbreak{}learning) \\
Classroom-based Experiments in Productive Failure & 2011 & OA & stöder påståendet (0.98) (supports\_\allowbreak{}claim; productive\_\allowbreak{}failure; classroom\_\allowbreak{}experiments; conceptual\_\allowbreak{}insight) \\
Designing for Productive Failure & 2011 & OA & stöder påståendet (0.97) (supports\_\allowbreak{}claim; productive\_\allowbreak{}failure; math; conceptual\_\allowbreak{}understanding; transfer) \\
Designing for productive failure in mathematical problem solving & 2009 & OA & stöder påståendet (0.97) (supports\_\allowbreak{}claim; productive\_\allowbreak{}failure; math; higher\_\allowbreak{}order\_\allowbreak{}application) \\
Differential benefits of explicit failure-driven and success-driven scaffolding in problem-solving prior to instruction. & 2021 & Scopus & stöder påståendet (0.89) (supports\_\allowbreak{}claim; failure\_\allowbreak{}driven\_\allowbreak{}scaffolding; productive\_\allowbreak{}failure; boundary\_\allowbreak{}conditions) \\
DIY productive failure: boosting performance in a large undergraduate biology course & 2019 & OA & stöder påståendet (0.97) (supports\_\allowbreak{}claim; productive\_\allowbreak{}failure; biology; large\_\allowbreak{}classroom) \\
Does sequence matter? Productive failure and designing online authentic learning for process engineering & 2016 & OA & stöder påståendet (0.97) (supports\_\allowbreak{}claim; productive\_\allowbreak{}failure; online\_\allowbreak{}learning; conceptual\_\allowbreak{}knowledge; transfer) \\
Flipping a simulation before instruction can improve students' learning, interest and perceived competence & 2026 & Scopus & stöder påståendet (0.79) (supports\_\allowbreak{}claim; flipping\_\allowbreak{}before\_\allowbreak{}instruction; simulation\_\allowbreak{}learning) \\
Learning by exploring: How much guidance is optimal? & 2019 & WoS & stöder påståendet (0.98) (qualifies\_\allowbreak{}claim; guidance\_\allowbreak{}during\_\allowbreak{}exploration; supports\_\allowbreak{}claim) \\
Learning from failure: a meta-analysis of the empirical studies & 2018 & WoS & stöder påståendet (0.93) (meta\_\allowbreak{}analysis; learning\_\allowbreak{}from\_\allowbreak{}failure; supports\_\allowbreak{}claim) \\
Motivation and transfer: The role of achievement goals in preparation for future learning & 2009 & OA & stöder påståendet (0.97) (supports\_\allowbreak{}claim; invention\_\allowbreak{}activities; transfer; motivation) \\
Motivation and Transfer: The Role of Mastery-Approach Goals in Preparation for Future Learning & 2012 & WoS & stöder påståendet (0.97) (supports\_\allowbreak{}claim; invention\_\allowbreak{}activities; preparation\_\allowbreak{}for\_\allowbreak{}future\_\allowbreak{}learning; motivation) \\
Observing a robot peer's failures facilitates students' classroom learning & 2025 & WoS & stöder påståendet (0.97) (productive\_\allowbreak{}failure; vicarious\_\allowbreak{}failure; supports\_\allowbreak{}claim; math) \\
Outcomes and Mechanisms of Transfer in Invention Activities & 2011 & OA & stöder påståendet (0.96) (supports\_\allowbreak{}claim; invention\_\allowbreak{}activities; adaptive\_\allowbreak{}knowledge; transfer) \\
Problem-solving prior to instruction in learning motor skills -- initial self-determined practice improves javelin throwing performance & 2022, 2023 & OA, WoS & stöder påståendet (0.96) (preprint och artikel av samma studie; supports\_\allowbreak{}claim; motor\_\allowbreak{}skills; self\_\allowbreak{}determined\_\allowbreak{}practice) \\
Productive Failure: A Hidden Efficacy of Seemingly Unproductive Production & 2006 & OA & stöder påståendet (0.98) (supports\_\allowbreak{}claim; productive\_\allowbreak{}failure; transfer; foundational\_\allowbreak{}study) \\
Reverse the routine: Problem solving before instruction improves conceptual knowledge in undergraduate physics & 2018 & Scopus & stöder påståendet (noterad, inte citerad; supports\_\allowbreak{}claim; physics; conceptual\_\allowbreak{}knowledge; pf\_\allowbreak{}unreachable) \\
Robust effects of the efficacy of explicit failure-driven scaffolding in problem-solving prior to instruction: A replication and extension & 2021 & OA & stöder påståendet (0.98) (supports\_\allowbreak{}claim; failure\_\allowbreak{}driven\_\allowbreak{}scaffolding; transfer; replication) \\
Wait for it . . . Delaying Instruction Improves Mathematics Problem Solving: A Classroom Study & 2014 & OA & stöder påståendet (qualifies\_\allowbreak{}claim; children; classroom\_\allowbreak{}study; knowledge\_\allowbreak{}application\_\allowbreak{}phase; supports\_\allowbreak{}claim; boundary\_\allowbreak{}young\_\allowbreak{}learners) \\
Advancing the Guidance Debate: Lessons from Educational Psychology and Implications for Biochemistry Learning & 2020 & OA & kvalificerar eller motsäger påståendet (0.97) (productive\_\allowbreak{}failure; guided\_\allowbreak{}inquiry; worked\_\allowbreak{}examples; comparative\_\allowbreak{}pedagogies; qualifies\_\allowbreak{}claim) \\
Can failure be made productive also in Bayesian reasoning? A conceptual replication study & 2024 & OA & kvalificerar eller motsäger påståendet (0.98) (qualifies\_\allowbreak{}claim; conceptual\_\allowbreak{}replication; bayesian\_\allowbreak{}reasoning; boundary\_\allowbreak{}conditions) \\
Collaborative or Individual Learning within Productive Failure: Does the Social Form of Learning Make a Difference? & 2015 & OA & kvalificerar eller motsäger påståendet (0.98) (productive\_\allowbreak{}failure; collaboration\_\allowbreak{}vs\_\allowbreak{}individual; qualifies\_\allowbreak{}claim; young\_\allowbreak{}learners) \\
Comparing the effectiveness of preparatory activities that help undergraduate students learn from instruction & 2022 & OA & kvalificerar eller motsäger påståendet (0.98) (preparatory\_\allowbreak{}activities; undergraduates; qualifies\_\allowbreak{}claim) \\
Conditions for Effective Learning Without Upfront Instruction: How Practice with Feedback Supports Memory, Generalization, Motivation, and Metacognition & 2026 & Scopus & kvalificerar eller motsäger påståendet (adjacent\_\allowbreak{}review; learning\_\allowbreak{}without\_\allowbreak{}upfront\_\allowbreak{}instruction; unclear\_\allowbreak{}ps\_\allowbreak{}i\_\allowbreak{}comparison; qualifies\_\allowbreak{}claim; boundary\_\allowbreak{}conditions; review\_\allowbreak{}or\_\allowbreak{}theory) \\
Discussing Student Solutions Is Germane for Learning when Providing or Delaying Instruction & 2015 & WoS & kvalificerar eller motsäger påståendet (0.94) (qualifies\_\allowbreak{}claim; instruction\_\allowbreak{}phase\_\allowbreak{}matters; comparing\_\allowbreak{}solutions) \\
Element interactivity as a factor influencing the effectiveness of worked example--problem solving and problem solving--worked example sequences & 2020 & Scopus & kvalificerar eller motsäger påståendet (0.90) (worked\_\allowbreak{}examples; problem\_\allowbreak{}solving\_\allowbreak{}sequence; element\_\allowbreak{}interactivity; qualifies\_\allowbreak{}claim) \\
Examining Productive Failure, Productive Success, Unproductive Failure, and Unproductive Success in Learning & 2016 & Scopus & kvalificerar eller motsäger påståendet (0.88) (productive\_\allowbreak{}failure; theory; qualifies\_\allowbreak{}claim) \\
Fail, flip, fix, feed: impact on engagement and learning outcomes & 2026 & WoS & kvalificerar eller motsäger påståendet (0.95) (productive\_\allowbreak{}failure; null\_\allowbreak{}or\_\allowbreak{}equivalent\_\allowbreak{}results; flipped\_\allowbreak{}classroom; qualifies\_\allowbreak{}claim) \\
How preparation-for-learning with a worked versus an open inventing problem affect subsequent learning processes in pre-service teachers & 2022 & OA & kvalificerar eller motsäger påståendet (0.97) (inventing\_\allowbreak{}vs\_\allowbreak{}worked\_\allowbreak{}problem; qualifies\_\allowbreak{}claim; preparation\_\allowbreak{}for\_\allowbreak{}learning) \\
Instructional sequences in science teaching: considering element interactivity when sequencing inquiry-based investigation activities and explicit instruction & 2024 & Scopus & kvalificerar eller motsäger påståendet (0.84) (instructional\_\allowbreak{}sequences; science\_\allowbreak{}teaching; element\_\allowbreak{}interactivity; qualifies\_\allowbreak{}claim) \\
Inventing a solution and studying a worked solution prepare differently for learning from direct instruction & 2015 & Scopus & kvalificerar eller motsäger påståendet (0.90) (inventing\_\allowbreak{}vs\_\allowbreak{}worked\_\allowbreak{}solution; preparation\_\allowbreak{}for\_\allowbreak{}learning; qualifies\_\allowbreak{}claim) \\
Inventing Prepares Learning Motivationally, but a Worked-out Solution Enhances Learning Outcomes & 2013 & OA & kvalificerar eller motsäger påståendet (0.97) (qualifies\_\allowbreak{}claim; worked\_\allowbreak{}solution\_\allowbreak{}better; inventing\_\allowbreak{}vs\_\allowbreak{}worked\_\allowbreak{}example) \\
Learning about multivariable causality with interactive simulations: exploration before instruction may hurt immediate gains but benefits transfer & 2025 & OA & kvalificerar eller motsäger påståendet (0.97) (qualifies\_\allowbreak{}claim; simulation\_\allowbreak{}learning; transfer; immediate\_\allowbreak{}vs\_\allowbreak{}transfer\_\allowbreak{}tradeoff) \\
Preparatory effects of problem solving versus studying examples prior to instruction & 2021 & OA & kvalificerar eller motsäger påståendet (0.97) (problem\_\allowbreak{}solving\_\allowbreak{}vs\_\allowbreak{}examples\_\allowbreak{}prior\_\allowbreak{}to\_\allowbreak{}instruction; preparatory\_\allowbreak{}effects; qualifies\_\allowbreak{}claim) \\
Probing boundary conditions of Productive Failure and analyzing the role of young students' collaboration & 2019 & WoS & kvalificerar eller motsäger påståendet (0.98) (qualifies\_\allowbreak{}claim; boundary\_\allowbreak{}conditions; young\_\allowbreak{}students; collaboration) \\
Productive versus vicarious failure: Do students need to fail themselves in order to learn? & 2022 & OA & kvalificerar eller motsäger påståendet (0.97) (vicarious\_\allowbreak{}failure; productive\_\allowbreak{}failure; comparison; qualifies\_\allowbreak{}claim) \\
Putting Students on the Path to Learning: The Case for Fully Guided Instruction. & 2012 & OA & kvalificerar eller motsäger påståendet (0.94) (qualifies\_\allowbreak{}claim; fully\_\allowbreak{}guided\_\allowbreak{}instruction; cognitive\_\allowbreak{}load\_\allowbreak{}critique) \\
The Effect of Direct Instruction versus Discovery Learning on the Understanding of Science Lessons by Second Grade Students & 2008 & OA & kvalificerar eller motsäger påståendet (0.90) (qualifies\_\allowbreak{}claim; direct\_\allowbreak{}instruction\_\allowbreak{}vs\_\allowbreak{}discovery; children; null\_\allowbreak{}results) \\
When Productive Failure Fails & 2019 & Scopus & kvalificerar eller motsäger påståendet (0.88) (qualifies\_\allowbreak{}claim; productive\_\allowbreak{}failure\_\allowbreak{}fails; boundary\_\allowbreak{}conditions) \\
When to inquire and when to instruct: Prior knowledge and element interactivity in medical learning & 2026 & Scopus & kvalificerar eller motsäger påståendet (0.79) (when\_\allowbreak{}to\_\allowbreak{}inquire\_\allowbreak{}and\_\allowbreak{}when\_\allowbreak{}to\_\allowbreak{}instruct; prior\_\allowbreak{}knowledge; element\_\allowbreak{}interactivity; qualifies\_\allowbreak{}claim) \\
\end{longtable}
\endgroup

